% Generated by scripts/assemble.py; corrections belong in source/corrections.json.
\examhead{24}{中国矿业大学 2026 年数学分析试卷}{28}
% p029-b01; source PDF page 29
\par\Needspace{5\baselineskip}
\blockanchor{p029-b01}
\majorquestion{一}选择题. 每题5分, 共25分.\par
% p029-b02; source PDF page 29
\blockanchor{p029-b02}
\question{1}当 $x \to -\infty$ 时, $f(x)$ 不趋于 $a$ 的定义为( ).  
A. $\exists M > 0$, 对 $\forall \varepsilon > 0$, $\forall x < -M$, $|f(x) - a| \geq \varepsilon$.  
B. $\exists M > 0$, 对 $\forall \varepsilon > 0$, $\exists x_0 < -M$, $|f(x_0) - a| \geq \varepsilon$.  
C. $\exists \varepsilon_0 > 0$, 对 $\forall M > 0$, $\exists x_0 < -M$, $|f(x) - a| < \varepsilon_0$.  
D. $\exists \varepsilon_0 > 0$, 对 $\forall M > 0$, $\exists x_0 < -M$, $|f(x) - a| \geq \varepsilon_0$.\par
% p029-b03; source PDF page 29
\blockanchor{p029-b03}
\question{2}$f(x)=\left\{\begin{aligned}x+2,&x\geq0,\\ x-2,&x<0.\end{aligned}\right.$  则  $f(x)$  在 x=0 处() A. 连续. B. 左连续但不右连续. C. 右连续但不左连续. D. 既不左连续也不右连续.\par
% p029-b04; source PDF page 29
\blockanchor{p029-b04}
\question{3}变限积分 $f(x) = \int_{0}^{x} te^{t} \, \mathrm{d}t$，则 $f(x) = (\quad)$ A. $xe^{x} + e^{x} - 1$. B. $xe^{x} + e^{x}$. C. $e^{x}$. D.
\question{0}\par
% p029-b05; source PDF page 29
\blockanchor{p029-b05}
\question{4}若级数  $\sum_{n=1}^{\infty}a_{n}$  条件收敛, 则  $\sum_{n=1}^{\infty}\left(a_{n}+\frac{1}{n^{2}}\right)$  ( )
A. 条件收敛. B. 绝对收敛. C. 发散. D. 可能条件收敛也可能绝对收敛.\par
% p029-b06; source PDF page 29
\blockanchor{p029-b06}
\question{5}积分  $\oint_{L}\frac{x\,dy-y\,dx}{x^{2}+y^{2}}=(\quad)$ ，其中 L 为  $x^{2}+2y^{2}=1$ ，取逆时针方向.
A.
\question{0}B.  $\pi$ . C.  $2\pi$ . D.  $-\pi$ .\par
% p029-b07; source PDF page 29
\par\Needspace{5\baselineskip}
\blockanchor{p029-b07}
\majorquestion{二}填空题. 每题 5 分, 共 20 分.\par
% p029-b08; source PDF page 29
\blockanchor{p029-b08}
\question{6}求  $f(x)=(x-1)x^{\frac{2}{3}}$  的极小值 \_\_\_\_.\par
% p029-b09; source PDF page 29
\blockanchor{p029-b09}
\question{7}求极限 $\lim_{n\to \infty}\left[(n^2 +1)^{\frac{1}{n}} - 1\right]\sin \frac{n\pi}{2} =$\par
% p029-b10; source PDF page 29
\par\Needspace{6\baselineskip}
\blockanchor{p029-b10}
\question{8}设函数 $f(x)$ 是以 $2\pi$ 为周期的连续函数, 若 $f(x)$ 在 $[-\pi, \pi]$ 上的傅里叶级数展开式为\par
\nopagebreak[4]
% p029-b11; source PDF page 29
\blockanchor{p029-b11}
\displayformula{\frac {a _ {0}}{2} + \sum_ {n = 1} ^ {\infty} \left(a _ {n} \cos n x + b _ {n} \sin n x\right).}
% p029-b12; source PDF page 29
\blockanchor{p029-b12}
求  $\frac{1}{\pi}\int_{-\pi}^{\pi}f^{2}(x)\mathrm{d}x=$  \_\_\_\_.\par
% p029-b13; source PDF page 29
\blockanchor{p029-b13}
\question{9}求  $\ln(1-x+x^{2})$  在 x=0 处的幂级数展开式 \_\_\_\_.\par
% p029-b14; source PDF page 29
\par\Needspace{5\baselineskip}
\blockanchor{p029-b14}
\majorquestion{三}解答题. 每题 15 分, 共 45 分.\par
% p029-b15; source PDF page 29
\blockanchor{p029-b15}
\question{10}设 $f(x)$ 在 $x = 0$ 的某邻域内有定义, 且 $\lim_{x \to 0} \frac{f(x)}{x^2} = 1$, 求 $\lim_{x \to 0} \left(1 + \frac{f(x)}{x}\right)^{\frac{1}{x}}$.\par
% p029-b16; source PDF page 29
\par\Needspace{6\baselineskip}
\blockanchor{p029-b16}
\question{11}求曲面积分\par
\nopagebreak[4]
% p029-b17; source PDF page 29
\blockanchor{p029-b17}
\displayformula{\oiint_ {\Sigma} y ^ {2} z \mathrm{d} x \mathrm{d} y + x z \mathrm{d} y \mathrm{d} z + x ^ {2} y \mathrm{d} z \mathrm{d} x.}
% p029-b18; source PDF page 29
\blockanchor{p029-b18}
其中曲面 $\Sigma$ 为 $z = x^{2} + y^{2}, x^{2} + y^{2} = 1, z = 0$ 所围区域的表面，取外侧.\par
% p029-b19; source PDF page 29
\blockanchor{p029-b19}
\question{12}$f(x,y) = \left\{ \begin{array}{ll}\frac{xy(x^2 - y^2)}{x^2 + y^2}, & (x,y)\neq (0,0),\\ 0, & (x,y) = (0,0). \end{array} \right.$ 求 $f_{xy}(0,0)$ 和 $f_{yx}(0,0)$\par
% p029-b20; source PDF page 29
\par\Needspace{5\baselineskip}
\blockanchor{p029-b20}
\majorquestion{四}证明题. 每题 15 分, 共 60 分.\par
% p029-b21; source PDF page 29
\blockanchor{p029-b21}
\question{13}证明: $\operatorname{sgn}(x)=\begin{cases}-1,&x<0;\\0,&x=0;\\1,&x>0.\end{cases}$ 在 $[-1,1]$ 上黎曼可积, 但在 $[-1,1]$ 上并没有原函数.\par
% p030-b00; source PDF page 30
\blockanchor{p030-b00}
\question{14}已知 $f(x)$ 在 $[a, b]$ 上连续, 证明: $f(x)$ 在 $[a, b]$ 上有最大值.\par
% p030-b01; source PDF page 30
\blockanchor{p030-b01}
\question{15}判断当 $\alpha$ 取何值时, 函数列 $f_{n}(x) = \frac{x(\ln n)^{\alpha}}{n^{x}}$ 在 $[0, +\infty)$ 上一致收敛?\par
% p030-b02; source PDF page 30
\blockanchor{p030-b02}
\question{16}解答如下问题:\par
% p030-b03; source PDF page 30
\blockanchor{p030-b03}
(1) 已知级数 $\sum_{n=1}^{\infty} a_n, \sum_{n=1}^{\infty} c_n$ 收敛, 且 $a_n \leq b_n \leq c_n$, 判断 $\sum_{n=1}^{\infty} b_n$ 是否收敛? 若收敛, 请证明, 若不收敛, 请给出反例.\par
% p030-b04; source PDF page 30
\blockanchor{p030-b04}
(2) 对于 $\sum_{n=1}^{\infty} \frac{(-1)^{n-1}}{n^{p+\frac{1}{n}}}$, 根据 $p$ 的值讨论该级数何时条件收敛, 何时绝对收敛?\par
